% Options for packages loaded elsewhere
% Options for packages loaded elsewhere
\PassOptionsToPackage{unicode}{hyperref}
\PassOptionsToPackage{hyphens}{url}
\PassOptionsToPackage{dvipsnames,svgnames,x11names}{xcolor}
%
\documentclass[
  a4paper,
]{article}
\usepackage{xcolor}
\usepackage{amsmath,amssymb}
\setcounter{secnumdepth}{5}
\usepackage{iftex}
\ifPDFTeX
  \usepackage[T1]{fontenc}
  \usepackage[utf8]{inputenc}
  \usepackage{textcomp} % provide euro and other symbols
\else % if luatex or xetex
  \usepackage{unicode-math} % this also loads fontspec
  \defaultfontfeatures{Scale=MatchLowercase}
  \defaultfontfeatures[\rmfamily]{Ligatures=TeX,Scale=1}
\fi
\usepackage{lmodern}
\ifPDFTeX\else
  % xetex/luatex font selection
\fi
% Use upquote if available, for straight quotes in verbatim environments
\IfFileExists{upquote.sty}{\usepackage{upquote}}{}
\IfFileExists{microtype.sty}{% use microtype if available
  \usepackage[]{microtype}
  \UseMicrotypeSet[protrusion]{basicmath} % disable protrusion for tt fonts
}{}
\usepackage{setspace}
\makeatletter
\@ifundefined{KOMAClassName}{% if non-KOMA class
  \IfFileExists{parskip.sty}{%
    \usepackage{parskip}
  }{% else
    \setlength{\parindent}{0pt}
    \setlength{\parskip}{6pt plus 2pt minus 1pt}}
}{% if KOMA class
  \KOMAoptions{parskip=half}}
\makeatother
% Make \paragraph and \subparagraph free-standing
\makeatletter
\ifx\paragraph\undefined\else
  \let\oldparagraph\paragraph
  \renewcommand{\paragraph}{
    \@ifstar
      \xxxParagraphStar
      \xxxParagraphNoStar
  }
  \newcommand{\xxxParagraphStar}[1]{\oldparagraph*{#1}\mbox{}}
  \newcommand{\xxxParagraphNoStar}[1]{\oldparagraph{#1}\mbox{}}
\fi
\ifx\subparagraph\undefined\else
  \let\oldsubparagraph\subparagraph
  \renewcommand{\subparagraph}{
    \@ifstar
      \xxxSubParagraphStar
      \xxxSubParagraphNoStar
  }
  \newcommand{\xxxSubParagraphStar}[1]{\oldsubparagraph*{#1}\mbox{}}
  \newcommand{\xxxSubParagraphNoStar}[1]{\oldsubparagraph{#1}\mbox{}}
\fi
\makeatother


\usepackage{longtable,booktabs,array}
\usepackage{calc} % for calculating minipage widths
% Correct order of tables after \paragraph or \subparagraph
\usepackage{etoolbox}
\makeatletter
\patchcmd\longtable{\par}{\if@noskipsec\mbox{}\fi\par}{}{}
\makeatother
% Allow footnotes in longtable head/foot
\IfFileExists{footnotehyper.sty}{\usepackage{footnotehyper}}{\usepackage{footnote}}
\makesavenoteenv{longtable}
\usepackage{graphicx}
\makeatletter
\newsavebox\pandoc@box
\newcommand*\pandocbounded[1]{% scales image to fit in text height/width
  \sbox\pandoc@box{#1}%
  \Gscale@div\@tempa{\textheight}{\dimexpr\ht\pandoc@box+\dp\pandoc@box\relax}%
  \Gscale@div\@tempb{\linewidth}{\wd\pandoc@box}%
  \ifdim\@tempb\p@<\@tempa\p@\let\@tempa\@tempb\fi% select the smaller of both
  \ifdim\@tempa\p@<\p@\scalebox{\@tempa}{\usebox\pandoc@box}%
  \else\usebox{\pandoc@box}%
  \fi%
}
% Set default figure placement to htbp
\def\fps@figure{htbp}
\makeatother





\setlength{\emergencystretch}{3em} % prevent overfull lines

\providecommand{\tightlist}{%
  \setlength{\itemsep}{0pt}\setlength{\parskip}{0pt}}



 


% --- Layout & Document Setup ---
\usepackage[margin=1in]{geometry}
\usepackage{setspace}
\usepackage{xcolor}
\usepackage{graphicx}
\usepackage{float}
\usepackage{tcolorbox}
\usepackage{wrapfig}
\usepackage{booktabs}
\usepackage{multirow}
\usepackage{multicol}
\usepackage{makecell}
\usepackage{siunitx}
\usepackage{tabularx}
\usepackage{algorithm}
\usepackage{algpseudocode}

% --- Mathematical Enhancements ---
\usepackage{amsmath, amssymb, amsfonts}
\setlength{\abovedisplayskip}{6pt}
\setlength{\belowdisplayskip}{6pt}

% --- Reference & Tracking Systems ---
\usepackage{hyperref}
\usepackage[backend=biber]{biblatex}
\usepackage{lineno}
% \linenumbers % Uncomment to enable line numbering for faculty review [cite: 2026-02-08]

% --- Consolidated Header/Footer Logic for H1 Collection ---
\usepackage{fancyhdr}
\pagestyle{fancy}
\fancyhf{} 

% Standard Page Configuration
\fancyhead[L,R]{\small\color{red}PREPRINT/WORKING PAPER}
\fancyhead[C]{\small\color{red}\url{https://krishpn.github.io}}
\fancyfoot[L,R]{\small\color{red}PREPRINT/WORKING PAPER}
\fancyfoot[C]{\thepage}
\renewcommand{\headrulewidth}{0.4pt}
\renewcommand{\footrulewidth}{0.4pt}

% First Page (Title Page) Style Override
\fancypagestyle{plain}{
  \fancyhf{}
  \fancyhead[L,R]{\small\color{red}PREPRINT/WORKING PAPER}
  \fancyhead[C]{%
    \small\color{red}\url{https://krishpn.github.io}
  }
  \fancyfoot[L,R]{\small\color{red}PREPRINT/WORKING PAPER}
  \fancyfoot[C]{\thepage}
}


% --- Custom Math Commands ---
\newcommand{\R}{\mathbb{R}}
\newcommand{\N}{\mathbb{N}}
\newcommand{\Z}{\mathbb{Z}}
\newcommand{\quotes}[1]{``#1''}
\makeatletter
\@ifpackageloaded{caption}{}{\usepackage{caption}}
\AtBeginDocument{%
\ifdefined\contentsname
  \renewcommand*\contentsname{Table of contents}
\else
  \newcommand\contentsname{Table of contents}
\fi
\ifdefined\listfigurename
  \renewcommand*\listfigurename{List of Figures}
\else
  \newcommand\listfigurename{List of Figures}
\fi
\ifdefined\listtablename
  \renewcommand*\listtablename{List of Tables}
\else
  \newcommand\listtablename{List of Tables}
\fi
\ifdefined\figurename
  \renewcommand*\figurename{Figure}
\else
  \newcommand\figurename{Figure}
\fi
\ifdefined\tablename
  \renewcommand*\tablename{Table}
\else
  \newcommand\tablename{Table}
\fi
}
\@ifpackageloaded{float}{}{\usepackage{float}}
\floatstyle{ruled}
\@ifundefined{c@chapter}{\newfloat{codelisting}{h}{lop}}{\newfloat{codelisting}{h}{lop}[chapter]}
\floatname{codelisting}{Listing}
\newcommand*\listoflistings{\listof{codelisting}{List of Listings}}
\makeatother
\makeatletter
\makeatother
\makeatletter
\@ifpackageloaded{caption}{}{\usepackage{caption}}
\@ifpackageloaded{subcaption}{}{\usepackage{subcaption}}
\makeatother
\usepackage{bookmark}
\IfFileExists{xurl.sty}{\usepackage{xurl}}{} % add URL line breaks if available
\urlstyle{same}
\hypersetup{
  pdftitle={Bayesian Network Reconstruction of Interbank Exposures (Nepal)},
  pdfauthor={Krishna Neupane},
  pdfkeywords={networks},
  colorlinks=true,
  linkcolor={blue},
  filecolor={Maroon},
  citecolor={Blue},
  urlcolor={Blue},
  pdfcreator={LaTeX via pandoc}}


\title{Bayesian Network Reconstruction of Interbank Exposures (Nepal)}
\author{Krishna Neupane}
\date{2026-10-22}
\begin{document}
\maketitle
\begin{abstract}
text text
\end{abstract}


\setstretch{1.5}
\section{Paper 1 --- Bayesian Network Reconstruction of Interbank
Exposures
(Nepal)}\label{paper-1-bayesian-network-reconstruction-of-interbank-exposures-nepal}

\textbf{Status:} design settled, not yet implemented. \textbf{Target
venue:} ICLR / NeurIPS / ICAIF, ≤7 pages.

\textbf{One-line description:} Reconstruct the latent bilateral
interbank exposure matrix \(W_t\) for Nepali BFIs from NRB bulletin
data, using an SBM-structured, covariate-informed generative model with
a differentiable regulatory-constraint projection layer, amortized by a
GNN/GRU inference network.

This supersedes the earlier version of
\texttt{feature\_model\_mapping.md}, which had drifted from the actual
design (it invented a generic ``Bayesian-net structural-constraint''
layer built from accounting identities --- that was never part of the
real design and is corrected below).

\subsection{0. Problem framing and contribution to
policy}\label{problem-framing-and-contribution-to-policy}

\textbf{Problem framing (what ``the network'' means here).} We do not
pre-specify the edges of the interbank network. Instead, following the
network reconstruction literature (e.g., Squartini et al., 2018), we
treat the bilateral exposure structure as latent and infer it from
publicly observable aggregates --- balance-sheet marginals and
bank-level covariates --- within a constrained statistical framework.
Concretely: a node is a bank; a directed edge \(i \to j\) is
\(W_{ij,t}\), the exposure of bank \(i\) to bank \(j\); an edge
``exists'' when \(W_{ij,t}>0\). Neither the edge set nor the weights are
observed --- only the row/column marginals and covariates are. This is a
problem of statistical inference under partial observability, not
analysis over a topology that is already known or self-reported.
(Revisit/tighten this paragraph later --- current form is a first pass.)

\textbf{Contribution to policy.}

\begin{itemize}
\item
  Zero New Reporting Burden: No central bank in a resource-constrained
  setting is going to mandate new bilateral exposure reporting purely
  for a research exercise --- that is a costly, politically slow
  regulatory change. The contribution here is that this doesn't require
  that change: a usable, continuously-updatable estimate of systemic
  interconnectedness is produced from data NRB (or any
  similarly-resourced central bank) already publishes every month, with
  zero new reporting burden. That is, by relying exclusively on
  aggregate publicly available data, the model requires zero new
  regulatory compliannce overhead.
\item
  Unlocking an answerable Supervisory Question: At present, there is no
  way to ask \emph{who is exposed to whom} or \emph{which banks are
  structurally systemic} --- not because the question doesn't matter,
  but because answering it was assumed to require data that isn't
  collected. \textbf{This method shows it doesn't.} It converts a
  previously unanswerable supervisory question into an answerable one,
  cheaply, and produces a reusable pipeline that other analyses can plug
  into. In summary, the \emph{data barrier} can be bypassed
  computationally. The method hence turns the blind spot in systemic
  risk monitoring into generalized framework providing where the future
  policy framework implemention can and should be aimed at.
\item
  Generalization: This also generalizes beyond Nepal: most
  emerging-market central banks publish the same kind of aggregate
  call-report data and face the identical constraint, so the method is a
  template for resource-constrained regulators broadly, not a
  Nepal-specific workaround.
\end{itemize}

\subsection{1. Why this is its own
paper}\label{why-this-is-its-own-paper}

Separated from the causal/policy question (Paper 2) because the two are
methodologically unrelated: this paper's contribution is a
reconstruction method with no causal-identification content; Paper 2's
contribution is a causal-identification design that \emph{consumes} this
paper's output as a fixed input. Bundling them would force a methods
paper to carry a policy narrative it doesn't need, and vice versa.

\subsection{2. Notation}\label{notation}

\begin{longtable}[]{@{}
  >{\raggedright\arraybackslash}p{(\linewidth - 2\tabcolsep) * \real{0.5000}}
  >{\raggedright\arraybackslash}p{(\linewidth - 2\tabcolsep) * \real{0.5000}}@{}}
\toprule\noalign{}
\begin{minipage}[b]{\linewidth}\raggedright
Symbol
\end{minipage} & \begin{minipage}[b]{\linewidth}\raggedright
Meaning
\end{minipage} \\
\midrule\noalign{}
\endhead
\bottomrule\noalign{}
\endlastfoot
\(i, j \in \{1,...,N\}\) & individual banking and financial
institutions \\
\(t \in \{1,...,T\}\) & periods (NRB bulletin months) \\
\(X_{i,t} \in \mathbb{R}^d\) & observed node covariates for bank \(i\)
at time \(t\) \\
\(r_{i,t}, c_{i,t} \ge 0\) & observed interbank lending / borrowing
marginals \\
\(K_{i,t}\) & core capitalof the bank, that defines the the regulatory
single-obligor cap) \\
\(W_t \in \mathbb{R}_{\ge 0}^{N \times N}\) & The true \textbf{latent}
bilateral interbank exposure matrix --- never observed, the
reconstruction target \\
\(z_i \in \{1,...,K\}\) & latent block/tier membership (e.g.~core
vs.~periphery) \\
\end{longtable}

\subsection{3. Generative model}\label{generative-model}

\begin{itemize}
\tightlist
\item
  Block membership: Banks are sorted into latent tiers (SBM priors) :
  \(z_i \sim \text{Categorical}(\pi)\) --- (Eq. 1).
\item
  Block-pair intensity: \(\Lambda \in \mathbb{R}_{\ge0}^{K\times K}\)
  --- the SBM core
\item
  Fitness/degree-correction from covariates: bank baseline activity
  scales with covariates: \(\phi_{i,t} = \exp(w^\top X_{i,t})\) --- (Eq.
  2)
\item
  Bilateral Latent edge intensity: the exposure magnigute is a fucntion
  of tier interaction and invidivudal fitness.
  \[\lambda_{ij,t} = \Lambda_{z_i z_j} \cdot \phi_{i,t}\phi_{j,t} \cdot \exp(h_{i,t}^\top h_{j,t}) \qquad \text{(Eq. 3)}\]
\item
  Zero-inflated edge weight (most bank pairs have \emph{no} exposure).
  The edges are modeled via hurle process - a Bernoulli gate determines
  connection existence, followed by a LogNormal distribution for
  positive exposure volumes:
  \[\Pr(\tilde W_{ij,t}=0) = \sigma(-\eta_{ij,t}), \qquad \tilde W_{ij,t} \mid \tilde W_{ij,t}>0 \sim \text{LogNormal}(\log\lambda_{ij,t}, \sigma^2) \qquad \text{(Eq. 4)}\]
\item
  Marginals are a \textbf{deterministic function} of \(W_t\):
  \[r_{i,t} = \sum_j W_{ij,t}, \qquad c_{j,t} = \sum_i W_{ij,t} \qquad \text{(Eq. 5)}\]
  This means the marginal constraint belongs in the \emph{architecture}
  (a hard projection), not the likelihood --- see §5. Rows and columns
  sums mach observed micro-prudential totals (\(r_{i,t}, c_{i,t}\)),
  enforced strictly via the Sinkhorn projection layer rather than soft
  likelihood penalties.
\end{itemize}

\subsection{4. Inference network (amortized, PyG GNN +
GRU)}\label{inference-network-amortized-pyg-gnn-gru}

\begin{itemize}
\item
  Design Principle: The encoder processes covariates/history only --- it
  cannot message-pass over \(W_t\) itself, since that would be circular
  (the whole point is that \(W_t\) is unobserved). This design principle
  is implemented because in the standard graph autoencodersv(GAEs), the
  adjacency matrix defined by bilateral or multilateral relationshhips
  is an input to encoder which in turn uses message-passing across
  \emph{known connections} to learn the node embeddings. In the current
  context, the true interbank (\(W_t\)) network is not known publicly -
  that is why the paper aims to \emph{reconstruct} interbank network
  from publicly available data. Because the encoder has an input
  \(W_t\), but the goal is to figure out the what \(W_t\) is, then
  modeling is trapped into logical circularity which is why this
  principle is adopted in order answer before it becomes burden in
  downstream task.
\item
  Temporal encoder:
  \(h_{i,t} = \text{GRU}_\theta(X_{i,t}, h_{i,t-1}) \qquad \text{(Eq. 6)}\),
  optionally refined by a GNN layer over an auxiliary \emph{observed}
  graph if one exists (e.g.~common-sector exposure --- see §8).
\item
  Amortized block-assignment posterior:
  \(q_\psi(z_i \mid X_{i,1:t}) = \text{Categorical}(\text{softmax}(\text{MLP}_\psi(h_{i,t}))) \qquad \text{(Eq. 7)}\)
\item
  Pairwise decoder produces \((\eta_{ij,t}, \lambda_{ij,t})\) from
  \((h_{i,t}, h_{j,t}, z_i, z_j)\) --- full amortized variational
  posterior \(q_\psi(W_t \mid X_{1:t}, r_t, c_t)\)
\item
  Reparameterized sampling: Generates log-normal reparam trick for the
  continuous weight, and Gumbel-softmax for the discrete existence gate.
\end{itemize}

\subsection{5. Constraint-satisfaction layer (the technical
core)}\label{constraint-satisfaction-layer-the-technical-core}

Differentiable \textbf{capacitated Sinkhorn / entropic-optimal transport
projection}:

\[W_t = \text{Sinkhorn}_{\text{cap}}(\tilde W_t;\; r_t, c_t, \kappa K_t) \qquad \text{(Eq. 8)}\]

Enforces, exactly and differentiably:

\begin{itemize}
\tightlist
\item
  \begin{enumerate}
  \def\labelenumi{\arabic{enumi}.}
  \tightlist
  \item
    row/column marginals match the observed \(r_{i,t}, c_{i,t}\)
  \end{enumerate}
\item
  \begin{enumerate}
  \def\labelenumi{\arabic{enumi}.}
  \setcounter{enumi}{1}
  \tightlist
  \item
    the regulatory single-obligor cap \(W_{ij,t} \le \kappa K_{i,t}\)
  \end{enumerate}
\end{itemize}

References: Cuturi (2013) entropic OT; Mena et al.~(2018)
Gumbel-Sinkhorn. Flagged explicitly as ``the one genuinely nontrivial
implementation piece'' in the design discussion --- this is the part
most likely to need real engineering effort, not the GNN/GRU scaffolding
around it.

\textbf{Correction vs.~the earlier \texttt{feature\_model\_mapping.md}:}
there is no ``Bayesian-net structural-constraint edge'' tying
\texttt{TOTAL\ Liabilities\ =\ TOTAL\ ASSETS} or \texttt{Net\ Profit}
into the graph. That was never part of this design. The accounting
identities in the raw NRB statements are a \textbf{data-quality
cross-check} on the extractor (do the reported totals reconcile?), not a
component of the generative model or the constraint layer.

\subsection{6. Training objective}\label{training-objective}

\[\mathcal{L} = -\mathbb{E}_q[\log p(W_t^{\text{true}})]_{\text{synthetic only}} + \text{KL}(q_\psi(z)\,\|\,p(z)) + \lambda_{\text{smooth}}\sum_t \|W_t - W_{t-1}\|_F^2 + \lambda_{\text{reg}}\|\Lambda, w\|^2\]

\begin{itemize}
\item
  \textbf{Synthetic Validation data} (known ground-truth \(W_t\) from an
  SBM generator): the first term is a real supervised reconstruction
  loss --- this is the main validation evidence. That, when training on
  synthetic data generated via the SBM (where the true networks
  \(W_t^{true}\) are known), the first term acts as a suprvised
  reconstruction loss, which in turn serves as the primary validation of
  the model's accuracy.
\item
  \textbf{Real data} (no \(W_t\) is ever observed): When applied to a
  real world central bank data where \(W_t\) is unobserved, the
  supervised reconstruction term drops out entirely, the first term
  drops out entirely. What's left is a regularized/MAP-style constrained
  fit, \emph{not} likelihood-based estimation. That is, in a standard
  statiscal or ML models, the model looks into observed ground trught
  answers (eg \(W_t^{true}\)). The model is trained by calculating the
  likelihood of model's guesses against the actual answers followed by
  parameter optimization to maximize the likelhood. However in this
  case, the true \(W_t^{true}\) does not exist in publicly available
  data, therefore a Maximum Likelhood calculation falls apart.
  Therefore, by implementing MAP instead of model fitting into known
  potentially right answers, the model aims to learn the most possible
  configuration of the network given the structural priors (eg latent
  tiers \(z_i\)) and regulatory rules.
\item
  Real World Evaluation: Evaluation on real data has to rely on indirect
  evidence: temporal stability of the recovered \(W_t\) sequence, and
  comparison against a classical baseline (§8).
\end{itemize}

\subsection{7. Design rationale}\label{design-rationale}

\begin{itemize}
\tightlist
\item
  \textbf{SBM over preferential-attachment/fitness models} as the sole
  structural backbone: core- periphery tiering in interbank markets is
  literally a latent-block story, so SBM is the natural fit. PA/fitness
  models (Caldarelli et al.) are cited in related work for contrast but
  not built.
\item
  \textbf{GRU, not a temporal-GNN} (EvolveGCN, TGN, etc.) for the
  temporal mechanism: deliberately the simpler choice, given the ``big
  model, small data'' risk --- roughly 50--80 institutions, at most
  \textasciitilde170 periods. A fancier temporal-GNN would be
  over-parameterized for this panel size.
\item
  \textbf{Mandatory baseline}: head-to-head comparison against
  Gandy--Veraart / RAS maximum-entropy network reconstruction, on both
  the synthetic SBM data (known ground truth) and the real panel.
  Explicit instruction from the design discussion: report honestly if
  the deep model doesn't beat the classical baseline --- don't bury a
  negative result.
\end{itemize}

\subsection{8. Variable inventory (what actually feeds the
model)}\label{variable-inventory-what-actually-feeds-the-model}

Net statement from the design discussion: \textbf{the reconstruction
model needs \textasciitilde5 covariates per node per period, plus the
two interbank marginals. Everything else in the NRB bulletins is either
a Paper 2 outcome variable or not usable given current data coverage.}

\subsubsection{8a. The two marginals (non-negotiable --- the
reconstruction
target)}\label{a.-the-two-marginals-non-negotiable-the-reconstruction-target}

The names below (\texttt{2.b\ Interbank\ Borrowing},
\texttt{INTERBANK\ LENDING\ (Marginal\ r\_i)}) were this doc's own
shorthand, not literal bulletin text, and the row-marginal name
additionally implied a standalone top-level item that doesn't exist.
Confirmed against the real extracted CSV
(\texttt{nrb\_balance\_sheet.csv}) on 2026-10-22 --- see §9.4:

\begin{longtable}[]{@{}
  >{\raggedright\arraybackslash}p{(\linewidth - 6\tabcolsep) * \real{0.2500}}
  >{\raggedright\arraybackslash}p{(\linewidth - 6\tabcolsep) * \real{0.2500}}
  >{\raggedright\arraybackslash}p{(\linewidth - 6\tabcolsep) * \real{0.2500}}
  >{\raggedright\arraybackslash}p{(\linewidth - 6\tabcolsep) * \real{0.2500}}@{}}
\toprule\noalign{}
\begin{minipage}[b]{\linewidth}\raggedright
Symbol
\end{minipage} & \begin{minipage}[b]{\linewidth}\raggedright
Real bulletin text
\end{minipage} & \begin{minipage}[b]{\linewidth}\raggedright
Parent section
\end{minipage} & \begin{minipage}[b]{\linewidth}\raggedright
Statement / side
\end{minipage} \\
\midrule\noalign{}
\endhead
\bottomrule\noalign{}
\endlastfoot
\(c_{i,t}\) (column marginal) & \texttt{b.\ Interbank\ Borrowing} &
\texttt{2\ BORROWINGS} & Balance Sheet, \textbf{Liabilities} \\
\(r_{i,t}\) (row marginal) & \texttt{a.\ Interbank\ Lending} &
\texttt{3\ SHARE\ \&\ OTHER\ INVESTMENT} & Balance Sheet,
\textbf{Assets} \\
\end{longtable}

Both are lettered sub-items, not top-level rows ---
\texttt{r\_\{i,t\}\$} in particular is a breakdown component of
\texttt{3\ SHARE\ \&\ OTHER\ INVESTMENT}'s own reported total, which
matters for the size formula below (§9.4). These are the actual payoff
variables for the whole method --- everything else exists to help infer
the matrix whose row/column sums these fix.

\subsubsection{\texorpdfstring{8b. Node covariates \(X_{i,t}\)
(block-membership / fitness
prior)}{8b. Node covariates X\_\{i,t\} (block-membership / fitness prior)}}\label{b.-node-covariates-x_it-block-membership-fitness-prior}

\begin{longtable}[]{@{}
  >{\raggedright\arraybackslash}p{(\linewidth - 6\tabcolsep) * \real{0.2500}}
  >{\raggedright\arraybackslash}p{(\linewidth - 6\tabcolsep) * \real{0.2500}}
  >{\raggedright\arraybackslash}p{(\linewidth - 6\tabcolsep) * \real{0.2500}}
  >{\raggedright\arraybackslash}p{(\linewidth - 6\tabcolsep) * \real{0.2500}}@{}}
\toprule\noalign{}
\begin{minipage}[b]{\linewidth}\raggedright
\#
\end{minipage} & \begin{minipage}[b]{\linewidth}\raggedright
Covariate
\end{minipage} & \begin{minipage}[b]{\linewidth}\raggedright
Built from
\end{minipage} & \begin{minipage}[b]{\linewidth}\raggedright
Role
\end{minipage} \\
\midrule\noalign{}
\endhead
\bottomrule\noalign{}
\endlastfoot
1 & Size & sum of the \textbf{10} real top-level asset rows:
\texttt{1\ LIQUID\ FUNDS}, \texttt{2\ INVESTMENT\ IN\ SECURITIES},
\texttt{3\ SHARE\ \&\ OTHER\ INVESTMENT},
\texttt{4\ LOANS\ \&\ ADVANCES}, \texttt{5\ BILL\ PURCHASED},
\texttt{6\ LOANS\ AGAINST\ COLLECTED\ BILLS}, \texttt{7\ FIXED\ ASSETS},
\texttt{8\ OTHER\ ASSETS}, \texttt{9\ EXPENSES\ NOT\ WRITTEN\ OFF},
\texttt{10\ NON\ BANKING\ ASSETS} (\textbf{not} the
\texttt{TOTAL\ ASSETS} label directly --- see §9.1). \(r_{i,t}\) is
\textbf{not} added as an 11th term --- per §8a/§9.4 it's a sub-item
already folded into \texttt{3\ SHARE\ \&\ OTHER\ INVESTMENT}'s reported
total, so adding it separately would double-count it & SBM block/fitness
prior, scale control \\
2 & Capital ratio & \texttt{1\ CAPITAL\ FUND} (Liabilities) ÷ size &
solvency/fitness term; also sets the obligor cap \(K_{i,t}\) \\
3 & Liquidity ratio & \texttt{1\ LIQUID\ FUNDS} (Assets) ÷ size &
liquidity position \\
4 & Loan-to-deposit pressure & \texttt{4\ LOANS\ \&\ ADVANCES} (Assets)
÷ \texttt{3\ DEPOSITS} (Liabilities) & funding-pressure signal \\
5 & NRB-dependence & \texttt{a.\ NRB} ÷ \texttt{2\ BORROWINGS} (both
under Liabilities, \texttt{a.\ NRB} nested under \texttt{2\ BORROWINGS};
if \texttt{2\ BORROWINGS} is missing for a row --- some early-era
filings never print the top-level subtotal --- the denominator is
reconstructed as the sum of its sub-items: \texttt{a.\ NRB},
\texttt{b.\ Interbank\ Borrowing},
\texttt{c.\ Foreign\ Banks\ and\ Fin.\ Ins.},
\texttt{d.\ Other\ Financial\ Ins.}, \texttt{e.\ Bonds\ and\ Securities}
--- confirmed real text per §9.4; the \texttt{f.}/\texttt{g.} sub-items
this doc previously assumed existed under \texttt{2\ BORROWINGS} have
\textbf{not} been confirmed in the real data yet) & central-bank
liquidity dependence --- useful for the zero-inflation gate (a bank
funded mainly by NRB is less likely to have interbank exposure) \\
\end{longtable}

Optional secondary signal, not in the core 5: ownership type
(state/private/foreign), if available --- would add homophily structure
to the block prior.

\subsubsection{8c. Auxiliary observed-graph candidate (optional GNN
refinement input,
§4)}\label{c.-auxiliary-observed-graph-candidate-optional-gnn-refinement-input-4}

\begin{longtable}[]{@{}
  >{\raggedright\arraybackslash}p{(\linewidth - 2\tabcolsep) * \real{0.5000}}
  >{\raggedright\arraybackslash}p{(\linewidth - 2\tabcolsep) * \real{0.5000}}@{}}
\toprule\noalign{}
\begin{minipage}[b]{\linewidth}\raggedright
Source
\end{minipage} & \begin{minipage}[b]{\linewidth}\raggedright
Role
\end{minipage} \\
\midrule\noalign{}
\endhead
\bottomrule\noalign{}
\endlastfoot
Statement of Credit, sector-wise breakdown (12 sectors + Others) &
shared sector-concentration = a correlated-shock signal independent of
the interbank marginals. Candidate as the auxiliary \emph{observed}
graph the GNN layer can refine over, since it doesn't require \(W_t\)
itself. \\
\end{longtable}

\subsubsection{8d. Explicitly excluded from Paper
1}\label{d.-explicitly-excluded-from-paper-1}

\begin{itemize}
\tightlist
\item
  \texttt{6\ Bills\ Payable}, \texttt{5\ Other\ Liabilities}
  (Liabilities, minor residuals) --- not used
\item
  \texttt{6/10\ Reconciliation\ A/c} (both sides) --- accounting plug,
  not economically meaningful
\item
  \texttt{8\ Expenses\ not\ Written\ off},
  \texttt{9\ Non\ Banking\ Assets} (Assets) --- legacy/minor, not used
\item
  Everything in Profit \& Loss --- these are Paper 2 outcome variables,
  not reconstruction inputs (see \texttt{paper2\_causal\_spillover.md})
\item
  \texttt{nrb\_statement\_of\_credit.csv} as a panel source --- only 2
  years of coverage in current data, can't support a 14-year panel; not
  worth cleaning further unless more years turn up uncollected
\end{itemize}

\subsubsection{8e. One-to-one equation-symbol ↔ raw-variable
map}\label{e.-one-to-one-equation-symbol-raw-variable-map}

Every symbol that appears in Eq. (1)--(8) is listed below with exactly
what real NRB field (if any) it is computed from. Symbols with no raw
field are latent or learned --- they're estimated by the model, not read
off the bulletin.

\begin{longtable}[]{@{}
  >{\raggedright\arraybackslash}p{(\linewidth - 8\tabcolsep) * \real{0.2000}}
  >{\raggedright\arraybackslash}p{(\linewidth - 8\tabcolsep) * \real{0.2000}}
  >{\raggedright\arraybackslash}p{(\linewidth - 8\tabcolsep) * \real{0.2000}}
  >{\raggedright\arraybackslash}p{(\linewidth - 8\tabcolsep) * \real{0.2000}}
  >{\raggedright\arraybackslash}p{(\linewidth - 8\tabcolsep) * \real{0.2000}}@{}}
\toprule\noalign{}
\begin{minipage}[b]{\linewidth}\raggedright
Symbol
\end{minipage} & \begin{minipage}[b]{\linewidth}\raggedright
Appears in
\end{minipage} & \begin{minipage}[b]{\linewidth}\raggedright
Raw NRB field
\end{minipage} & \begin{minipage}[b]{\linewidth}\raggedright
Statement / side
\end{minipage} & \begin{minipage}[b]{\linewidth}\raggedright
Transformation
\end{minipage} \\
\midrule\noalign{}
\endhead
\bottomrule\noalign{}
\endlastfoot
\(r_{i,t}\) & Eq. (5), (8) & \texttt{a.\ Interbank\ Lending} (sub-item
of \texttt{3\ SHARE\ \&\ OTHER\ INVESTMENT}) & Balance Sheet, Assets &
identity (reported value) \\
\(c_{i,t}\) & Eq. (5), (8) & \texttt{b.\ Interbank\ Borrowing} (sub-item
of \texttt{2\ BORROWINGS}) & Balance Sheet, Liabilities & identity
(reported value) \\
\(K_{i,t}\) & Eq. (8), \(\kappa K_{i,t}\) cap &
\texttt{1\ CAPITAL\ FUND} & Balance Sheet, Liabilities & identity
(reported value) --- note: raw value here, \emph{not} the ratio used in
\(X^{(2)}_{i,t}\) below \\
\(X^{(1)}_{i,t}\) & Eq. (2), component of \(X_{i,t}\) & sum of the 10
real top-level asset rows (listed in §8b) & Balance Sheet, Assets & sum
(not the \texttt{TOTAL\ ASSETS} label --- see §9.1; \(r_{i,t}\)
excluded, see §8b/§9.3) \\
\(X^{(2)}_{i,t}\) & Eq. (2) & \texttt{1\ CAPITAL\ FUND} & Balance Sheet,
Liabilities & ÷ \(X^{(1)}_{i,t}\) (capital ratio) \\
\(X^{(3)}_{i,t}\) & Eq. (2) & \texttt{1\ LIQUID\ FUNDS} & Balance Sheet,
Assets & ÷ \(X^{(1)}_{i,t}\) (liquidity ratio) \\
\(X^{(4)}_{i,t}\) & Eq. (2) & \texttt{4\ LOANS\ \&\ ADVANCES} ÷
\texttt{3\ DEPOSITS} & Balance Sheet, Assets ÷ Liabilities & ratio
across statement sides \\
\(X^{(5)}_{i,t}\) & Eq. (2) & \texttt{a.\ NRB} ÷ \texttt{2\ BORROWINGS}
& Balance Sheet, Liabilities (both) & ratio within Liabilities;
denominator falls back to the sum of \texttt{a.}--\texttt{e.} sub-items
(§8b) when \texttt{2\ BORROWINGS} is unreported \\
\(w\) & Eq. (2) & --- & --- & \textbf{learned} (fitness weight vector,
\(\dim = 5\), no raw field) \\
\(z_i\) & Eq. (1), (3), (7) & --- & --- & \textbf{latent} (inferred via
\(q_\psi\), Eq. 7); optionally informed by ownership type if used
(§8b) \\
\(\Lambda\) & Eq. (3) & --- & --- & \textbf{learned} (\(K\times K\)
block-pair intensity matrix) \\
\(h_{i,t}\) & Eq. (3), (6), (7) & \(X^{(1\text{-}5)}_{i,t}\) (all 5
covariates, fed in as the GRU input vector) & --- & \textbf{learned
embedding}, GRU output --- not read directly off any single field \\
\(\theta\) (GRU weights) & Eq. (6) & --- & --- & \textbf{learned} \\
\(\psi\) (posterior-net weights) & Eq. (7) & --- & --- &
\textbf{learned} \\
\(\eta_{ij,t}\) & Eq. (4) & --- & --- & \textbf{learned} (zero-inflation
logit, decoder output) \\
\(\sigma^2\) & Eq. (4) & --- & --- & \textbf{learned} (log-normal
dispersion --- global or per-block, §10) \\
\(\kappa\) & Eq. (8) & --- & --- & \textbf{exogenous regulatory
constant} (single-obligor exposure limit; not reported in the bulletin
itself --- a known regulatory parameter) \\
\(A^{\text{obs}}\) (auxiliary GNN graph, §4) & Eq. (6) refinement,
optional & sector-wise credit, 12 sectors & Statement of Credit &
co-exposure adjacency built from shared-sector lending shares \\
\end{longtable}

\textbf{Reading guide}:
\(X_{i,t} = \big(X^{(1)}_{i,t}, ..., X^{(5)}_{i,t}\big) \in \mathbb{R}^5\)
is the concrete instantiation of the generic \(X_{i,t}\in\mathbb{R}^d\)
in §2 --- i.e.~\(d=5\) for this model, and the ordering above (size,
capital ratio, liquidity ratio, LDR, NRB-dependence) is the fixed
component order used everywhere \(X_{i,t}\) appears in Eq. (2) and (6).

\subsection{9. Known data-pipeline caveats (apply before this goes into
the
model)}\label{known-data-pipeline-caveats-apply-before-this-goes-into-the-model}

\begin{enumerate}
\def\labelenumi{\arabic{enumi}.}
\tightlist
\item
  \textbf{\texttt{item\_code} collision}: confirmed by re-running
  \texttt{nrb\_extractor.py} against a real bulletin ---
  \texttt{item\_code} is not unique across the Balance Sheet's
  Liabilities and Assets sides. Both sides independently restart
  numbering at 1 (e.g.~\texttt{1\ CAPITAL\ FUND} on Liabilities
  vs.~\texttt{1\ LIQUID\ FUNDS} on Assets; same collision pattern for
  \texttt{2}, \texttt{3}, and lettered sub-items
  \texttt{a.}--\texttt{g.}). The extractor
  (\texttt{nrb\_extractor.py:197-269}, \texttt{\_extract\_rows})
  currently has no column marking which side a row is on --- that
  currently has to be inferred externally from which \texttt{TOTAL}
  block (\texttt{TOTAL\ \ \ \ Liabilities} vs.~\texttt{TOTAL\ ASSETS}) a
  row falls between on the page. \textbf{Fix needed}: add a
  \texttt{side} column to the extractor output before any join on
  \texttt{item\_code} alone, or it will silently collide liabilities and
  assets rows that share a code (e.g.~covariates 2 and 3 above, which
  both use an \texttt{item\_code\ ==\ 1} row from different sides).
\item
  Rotated bank-name header recovery (\texttt{recover\_bank\_names},
  \texttt{nrb\_extractor.py:130-171}) and the layout-mode row extraction
  are both working as validated against the real sample PDF --- no open
  bug there as of this pass.
\item
  \textbf{\texttt{item\_code} drift across eras}: the top-level code
  attached to a given labeled row is not stable over the panel's history
  --- e.g.~\texttt{LOANS\ \&\ ADVANCES} is coded \texttt{4} in some eras
  and \texttt{4.1} in others (with \texttt{BILL\ PURCHASED} shifting
  from \texttt{5} to \texttt{4.2} correspondingly). Matching must key
  off the normalized label text (and, for lettered sub-items, parent
  section), never off \texttt{item\_code} alone.
\item
  \textbf{Lettered sub-items (\texttt{a.}--\texttt{g.}) are reused
  across unrelated parent sections} (e.g.~\texttt{a.} means
  \texttt{Paid-up\ Capital} under \texttt{1\ CAPITAL\ FUND},
  \texttt{Cash\ Balance} under \texttt{1\ LIQUID\ FUNDS}, \texttt{NRB}
  under \texttt{2\ BORROWINGS}, and \texttt{Interbank\ Lending} under
  \texttt{3\ SHARE\ \&\ OTHER\ INVESTMENT} --- all in the same
  bulletin). The letter alone is not a key; a row's identity is
  (normalized label text, nearest preceding top-level/bare-number row in
  extraction order). The extractor should carry this parent link
  explicitly (a \texttt{parent\_item} or \texttt{section} column) rather
  than requiring every consumer to replay the state machine itself.
\item
  \textbf{Truncated-text duplicate rows}: some extractions clip a label
  mid-word (e.g.~\texttt{4\ LOANS\ \&} alongside the correct
  \texttt{4\ LOANS\ \&\ ADVANCES}, or \texttt{1\ CAPITAL} alongside
  \texttt{1\ CAPITAL\ FUND}) for the same \texttt{item\_code}/concept.
  Confirmed via \texttt{cut\ -d,\ -f7,8\ ...\ \textbar{}\ sort\ -u} on
  the real CSV (2026-10-22). Treat these as the same field (prefix match
  or longest-label-wins), not as missing data for the truncated variant.
\end{enumerate}

\subsection{10. Open items (not yet
decided)}\label{open-items-not-yet-decided}

\begin{itemize}
\tightlist
\item
  \textbf{\(K\) (number of SBM blocks)}: start with \(K=2\)
  (core/periphery), treat larger \(K\) as a robustness check, not the
  headline specification.
\item
  \textbf{\(\sigma^2\) (log-normal dispersion)}: learned, either
  globally or per-block --- not yet decided which.
\item
  \textbf{Venue/scope hook}: ICAIF fit is comfortable at the current
  scope (methods + real-data case study). ICLR/NeurIPS main track would
  need an added theoretical hook --- e.g.~a provable statement on how
  reconstruction uncertainty in \(\hat W_t\) propagates into coverage
  guarantees for a downstream interference estimator (which would also
  strengthen the Paper 1 → Paper 2 link). \textbf{Not resolved} ---
  needs a decision before drafting.
\end{itemize}

\subsection{11. Algorithms (pseudocode)}\label{algorithms-pseudocode}

Three algorithms, in dependency order: the capacitated Sinkhorn
projection (the hard constraint layer, Eq. 8) is standalone and testable
without the rest of the model; the generative forward pass (Eq. 1--5, 8)
samples a synthetic \(W_t\) for validation; training (Eq. 6--9) wraps
the amortized encoder/decoder of §4 around the same projection layer.
All three are written over \emph{symbol-level} inputs
(\(N, K, d, \kappa\), etc.) --- nothing below is a fixed dataset size or
constant; see §12 for how these get threaded through as config, not
literals.

\begin{algorithm}
\caption{\textsc{CapacitatedSinkhornProjection} (Eq. 8)}
\begin{algorithmic}[1]
\Require $\tilde{W} \in \mathbb{R}_{\ge 0}^{N \times N}$ (raw/sampled edge weights), $r, c \in \mathbb{R}_{\ge 0}^N$ (target marginals), $\text{cap} \in \mathbb{R}_{\ge 0}^N$ (per-row cap, $\text{cap}_i = \kappa \cdot K_{i,t}$), $\varepsilon > 0$ (entropic regularization), $n_{\text{iters}}, \text{tol}$ (stopping criteria)
\Ensure $W \in \mathbb{R}_{\ge 0}^{N \times N}$ satisfying row/col marginals and $W_{ij} \le \text{cap}_i$ exactly
\State $\text{Kern} \leftarrow \exp(\log(\tilde{W} + \delta) / \varepsilon)$ \Comment{$\delta$ tiny avoids $\log(0)$; config, not literal}
\State $u \leftarrow \mathbf{1}_N$; $v \leftarrow \mathbf{1}_N$
\For{$\text{iter} = 1 \dots n_{\text{iters}}$}
    \State $u_{\text{new}} \leftarrow r / (\text{Kern} \cdot v)$
    \State $u_{\text{new}} \leftarrow \min\left(u_{\text{new}}, \text{cap} / (\text{Kern} \cdot v)\right)$ \Comment{Enforce row-wise cap}
    \State $v_{\text{new}} \leftarrow c / (\text{Kern}^T \cdot u_{\text{new}})$
    \If{$\|u_{\text{new}} - u\| < \text{tol}$ \textbf{and} $\|v_{\text{new}} - v\| < \text{tol}$} \State \textbf{break}
    \EndIf
    \State $u, v \leftarrow u_{\text{new}}, v_{\text{new}}$
\EndFor
\State $W \leftarrow \text{diag}(u) \cdot \text{Kern} \cdot \text{diag}(v)$
\State \textbf{return} $W$
\end{algorithmic}
\end{algorithm}
\begin{algorithm}
\caption{\textsc{SampleSyntheticNetwork} (Eq. 1–5, 8; ground truth for validation, §6)}
\begin{algorithmic}[1]
\Require $N, K, d$ (problem dimensions), $\pi \in \Delta^K, \Lambda \in \mathbb{R}_{\ge 0}^{K \times K}, w \in \mathbb{R}^d, \sigma^2$ (generative params), $X_{1:N} \in \mathbb{R}^{N \times d}$ (covariates), $\kappa, K_{\text{cap}} \in \mathbb{R}^N$ (regulatory cap inputs)
\Ensure $W_{\text{true}} \in \mathbb{R}_{\ge 0}^{N \times N}, r, c \in \mathbb{R}^N$ (ground-truth network + marginals)
\For{$i = 1 \dots N$} 
    \State $z_i \sim \text{Categorical}(\pi)$ \Comment{Eq. 1}
\EndFor
\For{$i = 1 \dots N$} 
    \State $\varphi_i \leftarrow \exp(w^T X_i)$ \Comment{Eq. 2}
\EndFor
\For{$i, j = 1 \dots N$ ($i \neq j$)}
    \State $\lambda_{ij} \leftarrow \Lambda[z_i, z_j] \cdot \varphi_i \cdot \varphi_j$ \Comment{Eq. 3}
    \State $e_{ij} \sim \text{Bernoulli}(\sigma(\eta_{ij}))$ \Comment{Eq. 4, zero-inflation param}
    \State $\tilde{W}_{ij} \leftarrow e_{ij} \cdot \text{LogNormalSample}(\log \lambda_{ij}, \sigma^2)$
\EndFor
\State $r \leftarrow \text{row\_sums}(\tilde{W})$; $c \leftarrow \text{col\_sums}(\tilde{W})$
\State $W_{\text{true}} \leftarrow \textsc{CapacitatedSinkhornProjection}(\tilde{W}, r, c, \kappa \cdot K_{\text{cap}})$ \Comment{Eq. 8}
\State \textbf{return} $W_{\text{true}}, \text{row\_sums}(W_{\text{true}}), \text{col\_sums}(W_{\text{true}})$
\end{algorithmic}
\end{algorithm}
\begin{algorithm}
\caption{\textsc{TrainStep} (Eq. 6–9; amortized inference + one optimizer step)}
\begin{algorithmic}[1]
\Require Batch of periods $\{X_{1:N,t}, r_t, c_t, K_t\}_{t \in \text{batch}}$, optional $W_{\text{true},t}$, model params $\theta, \psi, \Lambda, w, \sigma^2$, hyperparams $\lambda_{\text{smooth}}, \lambda_{\text{reg}}, \kappa$
\Statex \textbf{State:} $h_{i,t-1}$ carried across periods (GRU hidden state)
\Ensure Loss value; updated $\theta, \psi, \Lambda, w$ after one gradient step
\For{$t \in \text{batch}$ (temporal order)}
    \For{$i = 1 \dots N$}
        \State $h_{i,t} \leftarrow \text{GRU}_\theta(X_{i,t}, h_{i,t-1})$ \Comment{Eq. 6}
        \State [Optional] $h_{i,t} \leftarrow \text{GNN\_refine}(h_{i,t}, A_{\text{obs},t})$ \Comment{§4, §8c}
        \State $q_i \leftarrow \text{softmax}(\text{MLP}_\psi(h_{i,t}))$ \Comment{Eq. 7}
        \State $z_i \sim \text{GumbelSoftmaxSample}(q_i)$ \Comment{Reparameterized}
    \EndFor
    \For{$i, j = 1 \dots N$ ($i \neq j$)}
        \State $(\eta_{ij,t}, \lambda_{ij,t}) \leftarrow \text{Decoder}_\psi(h_{i,t}, h_{j,t}, z_i, z_j)$ \Comment{Eq. 3–4}
        \State $e_{ij,t} \sim \text{GumbelSigmoidSample}(\eta_{ij,t})$ \Comment{Reparameterized gate}
        \State $\tilde{W}_{ij,t} \leftarrow e_{ij,t} \cdot \text{LogNormalSample}(\log \lambda_{ij,t}, \sigma^2)$
    \EndFor
    \State $W_t \leftarrow \textsc{CapacitatedSinkhornProjection}(\tilde{W}_t, r_t, c_t, \kappa \cdot K_t)$ \Comment{Eq. 8, Alg. 1}
    \If{$W_{\text{true},t}$ is given (synthetic branch)}
        \State $L_{\text{recon},t} \leftarrow -\log p(W_{\text{true},t} \mid \lambda_t, \eta_t, \sigma^2)$ \Comment{Eq. 9, first term}
    \Else
        \State $L_{\text{recon},t} \leftarrow 0$ \Comment{Real data: no ground truth}
    \EndIf
    \State $L_{\text{KL},t} \leftarrow \text{KL}(q_\psi(z_t) \parallel p(z))$
\EndFor
\State $L_{\text{smooth}} \leftarrow \sum_t \|W_t - W_{t-1}\|_F^2$ \Comment{Over consecutive periods in batch}
\State $L_{\text{reg}} \leftarrow \|\Lambda\|^2 + \|w\|^2$
\State $L \leftarrow \sum_t (L_{\text{recon},t} + L_{\text{KL},t}) + \lambda_{\text{smooth}} \cdot L_{\text{smooth}} + \lambda_{\text{reg}} \cdot L_{\text{reg}}$
\State Backpropagate $L$; update $\theta, \psi, \Lambda, w$ via optimizer step
\State \textbf{return} $L$
\end{algorithmic}
\end{algorithm}

Sinkhorn (operating on a non-negative matrix, enforcing row/col
marginals ≥ 0 via the kernel) is only the last step. The full pipeline
per Algorithm 3 (§11) and Eq. 1-8:

\begin{itemize}
\tightlist
\item
  \begin{enumerate}
  \def\labelenumi{\arabic{enumi}.}
  \tightlist
  \item
    SBM block assignment + covariate fitness → gives each bank pair a
    prior edge intensity \(\lambda_{ij,t}\) (Eq. 1-3), informed by
    \(z_i\) and the 5 covariates.
  \end{enumerate}
\item
  \begin{enumerate}
  \def\labelenumi{\arabic{enumi}.}
  \setcounter{enumi}{1}
  \tightlist
  \item
    GNN/GRU encoder + pairwise decoder → turns that into
    \((\eta_{ij,t}, \lambda_{ij,t})\) per pair, learned from data (Eq.
    6-7, §4).
  \end{enumerate}
\item
  \begin{enumerate}
  \def\labelenumi{\arabic{enumi}.}
  \setcounter{enumi}{2}
  \tightlist
  \item
    Zero-inflated log-normal sampling → produces the raw guess
    \(\tilde W_t\) (Eq.~4) --- this is where the actual ``who's
    connected to whom, roughly how much'' signal comes from.
  \end{enumerate}
\item
  \begin{enumerate}
  \def\labelenumi{\arabic{enumi}.}
  \setcounter{enumi}{3}
  \tightlist
  \item
    Only then does Sinkhorn project \(\tilde W_t\) onto the feasible set
    that exactly matches the observed marginals \(r_t,c_t\) and the
    regulatory cap (Eq.~8) --- a correction/projection step, not a
    generative one.
  \end{enumerate}
\end{itemize}

In summary, plain Sinkhorn/IPF applied to a uniform prior (no
covariates, no learned structure) is exactly the classical
Gandy-Veraart/RAS maximum-entropy reconstruction --- that's the baseline
this paper compares against (§7, §12 baseline requirement). The
contribution is that Sinkhorn here operates on a covariate- and
block-structure-informed \(\tilde W_t\), not blind uniform mass --- so
you justify why Sinkhorn

(it's the natural, differentiable way to enforce known hard constraints
exactly, à la Cuturi 2013) separately from why not Sinkhorn alone
(because unlike RAS/Gandy-Veraart, the prior it's projecting isn't
uninformed).

\section*{Methodological Literature Synthesis}

The model bridges three major streams of literature: classical financial
network reconstruction, entropic optimal transport, and amortized deep
latent variable modeling.

\subsection*{Classical Interbank Matrix Reconstruction}

Early efforts to estimate missing bilateral exposures (\(W_t\)) from
aggregated call reports relied primarily on Maximum Entropy (ME)
principles \cite{upper2011simulation}. While ME models successfully
match row and column balance sheet totals (\(r_t, c_t\)), they typically
yield fully connected, overly dense matrices that understate systemic
fragility. Modern maximum-entropy and fitness framework extensions
attempt to introduce sparsity \cite{squartini2018reconstruction}, yet
these classical models rely on rigid functional forms and cannot easily
learn complex temporal dynamics or non-linear node covariate signals
(\(X_{i,t}\)).

\subsection*{Entropic Optimal Transport \& Differentiable Constraints}

To enforce macro-prudential constraints differentiably, the model builds
on entropic optimal transport (EOT) frameworks.
\cite{cuturi2013sinkhorn} demonstrated that adding an entropic
regularizer to optimal transport transforms non-differentiable linear
programs into computationally efficient, GPU-parallelizable scaling
operations via the Sinkhorn-Knopp algorithm. Further extensions by
\cite{mena2018learning} showed that applying the Sinkhorn operator
directly to neural network outputs enables gradient backpropagation
through continuous matrix transformations and discrete topology
decisions. In our pipeline, this capacitated Sinkhorn projection
enforces row/column totals (\(r_t, c_t\)) and regulatory single-obligor
exposure limits (\(W_{ij,t} \le \kappa K_{i,t}\)) directly inside the
forward pass.

\subsection*{Amortized Variational Inference for Graphs}

Standard Graph Autoencoders require observing the graph adjacency matrix
to compute node representations. To bypass this circularity on
unobserved networks, our inference architecture draws inspiration from
Neural Relational Inference (NRI) \cite{kipf2018neural}. NRI established
that latent relational structures can be inferred purely from dynamic
node trajectories (\(X_{1:T}\)) using a graph variational autoencoder.
By conditioning the GNN/GRU temporal encoder strictly on independent
node covariates and history rather than unobserved edges, our pipeline
avoids circular identification logic while predicting continuous weight
distributions and discrete connection gates.

\subsection{8. Variable inventory (what actually feeds the
model)}\label{variable-inventory-what-actually-feeds-the-model-1}

Net statement from the design discussion: \textbf{the reconstruction
model needs \textasciitilde5 covariates per node per period, plus the
two interbank marginals. Everything else in the NRB bulletins is either
a Paper 2 outcome variable or not usable given current data coverage.}

\subsubsection{8a. The two marginals (non-negotiable --- the
reconstruction
target)}\label{a.-the-two-marginals-non-negotiable-the-reconstruction-target-1}

\begin{longtable}[]{@{}
  >{\raggedright\arraybackslash}p{(\linewidth - 4\tabcolsep) * \real{0.3333}}
  >{\raggedright\arraybackslash}p{(\linewidth - 4\tabcolsep) * \real{0.3333}}
  >{\raggedright\arraybackslash}p{(\linewidth - 4\tabcolsep) * \real{0.3333}}@{}}
\toprule\noalign{}
\begin{minipage}[b]{\linewidth}\raggedright
NRB field
\end{minipage} & \begin{minipage}[b]{\linewidth}\raggedright
Statement / side
\end{minipage} & \begin{minipage}[b]{\linewidth}\raggedright
Symbol
\end{minipage} \\
\midrule\noalign{}
\endhead
\bottomrule\noalign{}
\endlastfoot
\texttt{2.b\ Interbank\ Borrowing} & Balance Sheet, \textbf{Liabilities}
& \(c_{i,t}\) (column marginal) \\
\texttt{INTERBANK\ LENDING\ (Marginal\ r\_i)} & Balance Sheet,
\textbf{Assets} & \(r_{i,t}\) (row marginal) \\
\end{longtable}

These are the actual payoff variables for the whole method ---
everything else exists to help infer the matrix whose row/column sums
these fix.

\subsubsection{\texorpdfstring{8b. Node covariates \(X_{i,t}\)
(block-membership / fitness
prior)}{8b. Node covariates X\_\{i,t\} (block-membership / fitness prior)}}\label{b.-node-covariates-x_it-block-membership-fitness-prior-1}

\begin{longtable}[]{@{}
  >{\raggedright\arraybackslash}p{(\linewidth - 6\tabcolsep) * \real{0.2500}}
  >{\raggedright\arraybackslash}p{(\linewidth - 6\tabcolsep) * \real{0.2500}}
  >{\raggedright\arraybackslash}p{(\linewidth - 6\tabcolsep) * \real{0.2500}}
  >{\raggedright\arraybackslash}p{(\linewidth - 6\tabcolsep) * \real{0.2500}}@{}}
\toprule\noalign{}
\begin{minipage}[b]{\linewidth}\raggedright
\#
\end{minipage} & \begin{minipage}[b]{\linewidth}\raggedright
Covariate
\end{minipage} & \begin{minipage}[b]{\linewidth}\raggedright
Built from
\end{minipage} & \begin{minipage}[b]{\linewidth}\raggedright
Role
\end{minipage} \\
\midrule\noalign{}
\endhead
\bottomrule\noalign{}
\endlastfoot
1 & Size & sum of the 11 \texttt{ASSET\_TOP\_LEVEL} columns:
\texttt{1\ LIQUID\ FUNDS}, \texttt{2\ INVESTMENT\ IN\ SECURITIES},
\texttt{3\ SHARE\ \&\ OTHER\ INVESTMENT},
\texttt{4\ LOANS\ \&\ ADVANCES}, \texttt{5\ BILL\ PURCHASED},
\texttt{6\ LOANS\ AGAINST\ COLLECTED\ BILLS}, \texttt{7\ FIXED\ ASSETS},
\texttt{8\ OTHER\ ASSETS}, \texttt{9\ EXPENSES\ NOT\ WRITTEN\ OFF},
\texttt{10\ NON\ BANKING\ ASSETS},
\texttt{INTERBANK\ LENDING\ (Marginal\ r\_i)} (\textbf{not} the
\texttt{TOTAL\ ASSETS} label directly --- see data caveat below) & SBM
block/fitness prior, scale control \\
2 & Capital ratio & \texttt{1\ CAPITAL\ FUND} (Liabilities) ÷ size &
solvency/fitness term; also sets the obligor cap \(K_{i,t}\) \\
3 & Liquidity ratio & \texttt{1\ LIQUID\ FUNDS} (Assets) ÷ size &
liquidity position \\
4 & Loan-to-deposit pressure & \texttt{4\ LOANS\ \&\ ADVANCES} (Assets)
÷ \texttt{3\ DEPOSITS} (Liabilities) & funding-pressure signal \\
5 & NRB-dependence & \texttt{2.a\ Borrowings\ -\ NRB} ÷
\texttt{2\ BORROWINGS} (both Liabilities; if \texttt{2\ BORROWINGS} is
missing for a row --- some early-era filings never print the top-level
subtotal --- the denominator is reconstructed as the sum of its
sub-items: \texttt{2.a\ Borrowings\ -\ NRB},
\texttt{2.b\ Interbank\ Borrowing},
\texttt{2.c\ Foreign\ Banks\ and\ Financial\ Institutions},
\texttt{2.d\ Other\ Financial\ Institutions},
\texttt{2.e\ Bonds\ and\ Securities\ /\ NRB\ Bonds},
\texttt{2.f\ Borrowings\ -\ Govt\ Non-Financial\ Institutions},
\texttt{2.g\ Borrowings\ -\ Other\ Non-Financial\ Institutions}) &
central-bank liquidity dependence --- useful for the zero-inflation gate
(a bank funded mainly by NRB is less likely to have interbank
exposure) \\
\end{longtable}

Optional secondary signal, not in the core 5: ownership type
(state/private/foreign), if available --- would add homophily structure
to the block prior.

\subsubsection{8c. Auxiliary observed-graph candidate (optional GNN
refinement input,
§4)}\label{c.-auxiliary-observed-graph-candidate-optional-gnn-refinement-input-4-1}

\begin{longtable}[]{@{}
  >{\raggedright\arraybackslash}p{(\linewidth - 2\tabcolsep) * \real{0.5000}}
  >{\raggedright\arraybackslash}p{(\linewidth - 2\tabcolsep) * \real{0.5000}}@{}}
\toprule\noalign{}
\begin{minipage}[b]{\linewidth}\raggedright
Source
\end{minipage} & \begin{minipage}[b]{\linewidth}\raggedright
Role
\end{minipage} \\
\midrule\noalign{}
\endhead
\bottomrule\noalign{}
\endlastfoot
Statement of Credit, sector-wise breakdown (12 sectors + Others) &
shared sector-concentration = a correlated-shock signal independent of
the interbank marginals. Candidate as the auxiliary \emph{observed}
graph the GNN layer can refine over, since it doesn't require \(W_t\)
itself. \\
\end{longtable}

\subsubsection{8d. Explicitly excluded from Paper
1}\label{d.-explicitly-excluded-from-paper-1-1}

\begin{itemize}
\tightlist
\item
  \texttt{6\ Bills\ Payable}, \texttt{5\ Other\ Liabilities}
  (Liabilities, minor residuals) --- not used
\item
  \texttt{6/10\ Reconciliation\ A/c} (both sides) --- accounting plug,
  not economically meaningful
\item
  \texttt{8\ Expenses\ not\ Written\ off},
  \texttt{9\ Non\ Banking\ Assets} (Assets) --- legacy/minor, not used
\item
  Everything in Profit \& Loss --- these are Paper 2 outcome variables,
  not reconstruction inputs (see \texttt{paper2\_causal\_spillover.md})
\item
  \texttt{nrb\_statement\_of\_credit.csv} as a panel source --- only 2
  years of coverage in current data, can't support a 14-year panel; not
  worth cleaning further unless more years turn up uncollected
\end{itemize}

\subsubsection{8e. One-to-one equation-symbol ↔ raw-variable
map}\label{e.-one-to-one-equation-symbol-raw-variable-map-1}

Every symbol that appears in Eq. (1)--(8) is listed below with exactly
what real NRB field (if any) it is computed from. Symbols with no raw
field are latent or learned --- they're estimated by the model, not read
off the bulletin.

\begin{longtable}[]{@{}
  >{\raggedright\arraybackslash}p{(\linewidth - 8\tabcolsep) * \real{0.2000}}
  >{\raggedright\arraybackslash}p{(\linewidth - 8\tabcolsep) * \real{0.2000}}
  >{\raggedright\arraybackslash}p{(\linewidth - 8\tabcolsep) * \real{0.2000}}
  >{\raggedright\arraybackslash}p{(\linewidth - 8\tabcolsep) * \real{0.2000}}
  >{\raggedright\arraybackslash}p{(\linewidth - 8\tabcolsep) * \real{0.2000}}@{}}
\toprule\noalign{}
\begin{minipage}[b]{\linewidth}\raggedright
Symbol
\end{minipage} & \begin{minipage}[b]{\linewidth}\raggedright
Appears in
\end{minipage} & \begin{minipage}[b]{\linewidth}\raggedright
Raw NRB field
\end{minipage} & \begin{minipage}[b]{\linewidth}\raggedright
Statement / side
\end{minipage} & \begin{minipage}[b]{\linewidth}\raggedright
Transformation
\end{minipage} \\
\midrule\noalign{}
\endhead
\bottomrule\noalign{}
\endlastfoot
\(r_{i,t}\) & Eq. (5), (8) &
\texttt{INTERBANK\ LENDING\ (Marginal\ r\_i)} & Balance Sheet, Assets &
identity (reported value) \\
\(c_{i,t}\) & Eq. (5), (8) & \texttt{2.b\ Interbank\ Borrowing} &
Balance Sheet, Liabilities & identity (reported value) \\
\(K_{i,t}\) & Eq. (8), \(\kappa K_{i,t}\) cap &
\texttt{1\ CAPITAL\ FUND} & Balance Sheet, Liabilities & identity
(reported value) --- note: raw value here, \emph{not} the ratio used in
\(X^{(2)}_{i,t}\) below \\
\(X^{(1)}_{i,t}\) & Eq. (2), component of \(X_{i,t}\) & sum of the 11
\texttt{ASSET\_TOP\_LEVEL} columns (listed in §8b) & Balance Sheet,
Assets & sum (not the \texttt{TOTAL\ ASSETS} label --- see §9.1) \\
\(X^{(2)}_{i,t}\) & Eq. (2) & \texttt{1\ CAPITAL\ FUND} & Balance Sheet,
Liabilities & ÷ \(X^{(1)}_{i,t}\) (capital ratio) \\
\(X^{(3)}_{i,t}\) & Eq. (2) & \texttt{1\ LIQUID\ FUNDS} & Balance Sheet,
Assets & ÷ \(X^{(1)}_{i,t}\) (liquidity ratio) \\
\(X^{(4)}_{i,t}\) & Eq. (2) & \texttt{4\ LOANS\ \&\ ADVANCES} ÷
\texttt{3\ DEPOSITS} & Balance Sheet, Assets ÷ Liabilities & ratio
across statement sides \\
\(X^{(5)}_{i,t}\) & Eq. (2) & \texttt{2.a\ Borrowings\ -\ NRB} ÷
\texttt{2\ BORROWINGS} & Balance Sheet, Liabilities (both) & ratio
within Liabilities; denominator falls back to the sum of
\texttt{2.a}--\texttt{2.g} sub-items when \texttt{2\ BORROWINGS} is
unreported (§8b) \\
\(w\) & Eq. (2) & --- & --- & \textbf{learned} (fitness weight vector,
\(\dim = 5\), no raw field) \\
\(z_i\) & Eq. (1), (3), (7) & --- & --- & \textbf{latent} (inferred via
\(q_\psi\), Eq. 7); optionally informed by ownership type if used
(§8b) \\
\(\Lambda\) & Eq. (3) & --- & --- & \textbf{learned} (\(K\times K\)
block-pair intensity matrix) \\
\(h_{i,t}\) & Eq. (3), (6), (7) & \(X^{(1\text{-}5)}_{i,t}\) (all 5
covariates, fed in as the GRU input vector) & --- & \textbf{learned
embedding}, GRU output --- not read directly off any single field \\
\(\theta\) (GRU weights) & Eq. (6) & --- & --- & \textbf{learned} \\
\(\psi\) (posterior-net weights) & Eq. (7) & --- & --- &
\textbf{learned} \\
\(\eta_{ij,t}\) & Eq. (4) & --- & --- & \textbf{learned} (zero-inflation
logit, decoder output) \\
\(\sigma^2\) & Eq. (4) & --- & --- & \textbf{learned} (log-normal
dispersion --- global or per-block, §10) \\
\(\kappa\) & Eq. (8) & --- & --- & \textbf{exogenous regulatory
constant} (single-obligor exposure limit; not reported in the bulletin
itself --- a known regulatory parameter) \\
\(A^{\text{obs}}\) (auxiliary GNN graph, §4) & Eq. (6) refinement,
optional & sector-wise credit, 12 sectors & Statement of Credit &
co-exposure adjacency built from shared-sector lending shares \\
\end{longtable}

\textbf{Reading guide}:
\(X_{i,t} = \big(X^{(1)}_{i,t}, ..., X^{(5)}_{i,t}\big) \in \mathbb{R}^5\)
is the concrete instantiation of the generic \(X_{i,t}\in\mathbb{R}^d\)
in §2 --- i.e.~\(d=5\) for this model, and the ordering above (size,
capital ratio, liquidity ratio, LDR, NRB-dependence) is the fixed
component order used everywhere \(X_{i,t}\) appears in Eq. (2) and (6).

\subsection{9. Known data-pipeline caveats (apply before this goes into
the
model)}\label{known-data-pipeline-caveats-apply-before-this-goes-into-the-model-1}

\begin{enumerate}
\def\labelenumi{\arabic{enumi}.}
\tightlist
\item
  \textbf{\texttt{item\_code} collision}: confirmed by re-running
  \texttt{nrb\_extractor.py} against a real bulletin ---
  \texttt{item\_code} is not unique across the Balance Sheet's
  Liabilities and Assets sides. Both sides independently restart
  numbering at 1 (e.g.~\texttt{1\ CAPITAL\ FUND} on Liabilities
  vs.~\texttt{1\ LIQUID\ FUNDS} on Assets; same collision pattern for
  \texttt{2}, \texttt{3}, and lettered sub-items
  \texttt{a.}--\texttt{g.}). The extractor
  (\texttt{nrb\_extractor.py:197-269}, \texttt{\_extract\_rows})
  currently has no column marking which side a row is on --- that
  currently has to be inferred externally from which \texttt{TOTAL}
  block (\texttt{TOTAL\ \ \ \ Liabilities} vs.~\texttt{TOTAL\ ASSETS}) a
  row falls between on the page. \textbf{Fix needed}: add a
  \texttt{side} column to the extractor output before any join on
  \texttt{item\_code} alone, or it will silently collide liabilities and
  assets rows that share a code (e.g.~covariates 2 and 3 above, which
  both use an \texttt{item\_code\ ==\ 1} row from different sides).
\item
  Rotated bank-name header recovery (\texttt{recover\_bank\_names},
  \texttt{nrb\_extractor.py:130-171}) and the layout-mode row extraction
  are both working as validated against the real sample PDF --- no open
  bug there as of this pass.
\end{enumerate}

\subsection{10. Open items (not yet
decided)}\label{open-items-not-yet-decided-1}

\begin{itemize}
\tightlist
\item
  \(K\) (number of SBM blocks): start with \(K=2\) (core/periphery),
  treat larger \(K\) as a robustness check, not the headline
  specification.
\item
  \(\sigma^2\) (log-normal dispersion): learned, either globally or
  per-block --- not yet decided which.
\item
  \textbf{Venue/scope hook}: ICAIF fit is comfortable at the current
  scope (methods + real-data case study). ICLR/NeurIPS main track would
  need an added theoretical hook --- e.g.~a provable statement on how
  reconstruction uncertainty in \(\hat W_t\) propagates into coverage
  guarantees for a downstream interference estimator (which would also
  strengthen the Paper 1 → Paper 2 link). \textbf{Not resolved} ---
  needs a decision before drafting.
\end{itemize}

\subsection{11. Paper structure (≤7
pages)}\label{paper-structure-7-pages}

\begin{longtable}[]{@{}
  >{\raggedright\arraybackslash}p{(\linewidth - 4\tabcolsep) * \real{0.3333}}
  >{\raggedright\arraybackslash}p{(\linewidth - 4\tabcolsep) * \real{0.3333}}
  >{\raggedright\arraybackslash}p{(\linewidth - 4\tabcolsep) * \real{0.3333}}@{}}
\toprule\noalign{}
\begin{minipage}[b]{\linewidth}\raggedright
Section
\end{minipage} & \begin{minipage}[b]{\linewidth}\raggedright
Length
\end{minipage} & \begin{minipage}[b]{\linewidth}\raggedright
Content
\end{minipage} \\
\midrule\noalign{}
\endhead
\bottomrule\noalign{}
\endlastfoot
Intro + related work & \textasciitilde1.5p & position against
Aronow-Samii/Leung (network assumed \emph{known}), Elliott-Golub-Jackson
(network assumed known, no reconstruction), Gandy-Veraart
(reconstruction without covariate-informed priors or regulatory
constraints) \\
Method & \textasciitilde2.5p & §3--6 above \\
Synthetic validation & \textasciitilde1.5p & SBM-generated ground truth,
baseline comparison \\
Real-data case study & \textasciitilde1p & Nepal panel,
temporal-stability evidence \\
Discussion/limitations & \textasciitilde0.5p & honest treatment of the
no-ground-truth-on-real-data limitation \\
\end{longtable}




\end{document}
